\documentclass[11pt]{article}
\usepackage[a4paper,margin=1in]{geometry}
\usepackage[T1]{fontenc}
\usepackage{amsmath,amssymb,booktabs,tabularx}
\usepackage[authoryear,round]{natbib}
\usepackage[hidelinks]{hyperref}
\usepackage{siunitx}
\sisetup{group-digits=integer,group-separator={,},group-minimum-digits=4}
\newcommand{\conc}{\mathrm{C}}
\newcommand{\spr}{\mathrm{S}}
\title{Patient Shortage in Histopathology:\\Classifier or Coverage Limitation?}
\author{Yannik Queisler}
\date{September 9, 2026}

\begin{document}
\raggedbottom
\setlength{\emergencystretch}{2em}
\maketitle
\begin{abstract}
\noindent Reducing the number of training patients at fixed class patch counts damages histopathology classification with frozen features, and the tested mitigation methods recover little of that loss. This experiment asks whether a different classifier can exploit class information that remains accessible under patient shortage, or whether broader patient coverage benefits several distinct readouts of the same representation. It compares the original multilayer perceptron, regularized multinomial logistic regression, and cosine nearest neighbours on the existing balanced concentrated and balanced spread allocations of BRACS and TCGA-UT. The comparison uses three existing patient-disjoint splits and holds the training patches fixed across classifiers within each support condition. Patient-macro balanced accuracy measures classifier gains, coverage benefits, and their interaction. Classifier gains identify predictive structure that the MLP leaves unused, while consistent coverage benefits support the value of broader training-patient diversity.
\end{abstract}

\tableofcontents
\clearpage

\section{Motivation}
\label{sec:question}
The preceding benchmark found that independent-support deprivation reduced balanced accuracy by \num{2.62}--\num{3.03} percentage points on BRACS and \num{4.76}--\num{4.83} points on TCGA-UT, despite matched class patch counts. No exploratory method produced a positive BRACS recovery interval excluding zero; the strongest TCGA-UT recoveries closed only about one sixth of the deficit \citep{queisler2026benchmarkResults}. These findings establish damage that the tested interventions largely failed to repair. They do not establish that repair is impossible.

The experiment distinguishes two explanations for this loss. Under a \emph{classifier limitation}, the frozen features still contain useful predictive structure, but the original multilayer perceptron (MLP) does not exploit it well with the available training data. A different classifier could therefore improve accuracy using exactly the same shortage allocation. Under a \emph{coverage limitation}, the training set contains too few distinct patients to represent the variation encountered in unseen patients. Broader patient coverage should then benefit several classifiers, even if changing the classifier alone offers little improvement. Both limitations can contribute.

The encoder maps each patch \(x\) to a fixed feature vector \(z=\operatorname{Virchow2}(x)\). Reducing the training-patient set changes the labelled examples available to learn a classifier; the features of a given test patch remain fixed. The central question is therefore whether better decisions can be learned from the shortage data, and how much additional patient coverage helps. This report calls this ability \emph{classification generalization to unseen patients}: how successfully a classifier trained on the available frozen features predicts classes in held-out patients.

Two alternative \emph{readouts}, or classifiers applied to the frozen features, complement the original MLP. A linear probe tests whether a linear decision rule can use the available class structure \citep{alain2016probes}. A nearest-neighbour classifier tests whether nearby feature vectors have informative labels; this also provides an established way to evaluate frozen representations \citep{caron2021emerging}. Their gains under shortage and their responses to broader coverage distinguish the two explanations within the tested classifier families.
\section{Controlled comparison}
\label{sec:design}
The experiment retains BRACS's seven lesion classes and TCGA-UT's 30 eligible cancer types, including the benchmark's exclusions and label order. Each patch retains its frozen 2560-dimensional Virchow2 feature vector, formed by concatenating the CLS token and mean patch token. All eligibility rules and the three existing train--validation--test splits remain fixed \citep{queisler2026benchmarkProtocol}. A partitioning unit is a BRACS patient or TCGA-UT participant; both are called patients below. Patients are disjoint between training, validation, and test within each split. Repeated splits can contain the same patient in different roles.

The comparison uses balanced class allocations. Concentrated support (\(\conc\)) draws the available patches from fewer patients; spread support (\(\spr\)) uses broader training-patient coverage with the same class-specific patch counts. Both conditions retain the benchmark's exact patch identities: the shared concentrated-balanced allocation and the spread allocation under the native class assignment.

\begin{table}[htbp]
\centering
\small
\renewcommand{\arraystretch}{1.15}
\begin{tabularx}{\textwidth}{@{}p{0.22\textwidth}XX@{}}
\toprule
Design element & Fixed or varied & Purpose \\
\midrule
Datasets and splits & BRACS and TCGA-UT; three existing splits each & Evaluate both tasks in their existing cohorts. \\
Patient support & Balanced concentrated and balanced spread & Change patient breadth at matched class patch counts. \\
Readouts & Original CE MLP, logistic regression, cosine \(k\)-NN & Compare learned nonlinear, linear, and local decisions. \\
Evidence & Same frozen features; each allocation's exact patches & Give every readout the same labelled training evidence. \\
Evaluation & Unchanged natural validation and test patients & Measure generalization to unseen patients within each split. \\
\bottomrule
\end{tabularx}
\caption{Patient coverage and classifier choice are varied while each support allocation supplies identical training evidence to all readouts.}
\label{tab:design}
\end{table}

\noindent The historical comparison retains \num{13982}--\num{27202} training patches per BRACS split, with \num{78}--\num{89} concentrated patients versus \num{103} spread patients. TCGA-UT retains \num{36988}--\num{52204} patches, with \num{671}--\num{917} concentrated patients versus \num{4983} spread patients \citep{queisler2026benchmarkResults}. Different shortage doses prevent interpreting a larger effect on one dataset as intrinsic dataset susceptibility.

The spread condition supplies additional patient identities from the original training partition and provides a reference for performance with broader coverage. It changes patient identities and within-patient patch contributions together, so its effect concerns the full coverage intervention, including patient composition and influence.
\section{Readouts}
\label{sec:readouts}
Every readout receives the same training patches within each support condition. Model fitting uses training patients, hyperparameter selection uses validation patients, and final evaluation uses test patients.

\subsection{Original classifier}
The reference is the benchmark's two-layer MLP, trained with cross-entropy loss: a 512-unit hidden layer, ReLU activation, dropout, and a linear class output. The comparison uses the five existing confirmation initializations and their selected configurations, matched to the same features, labels, allocations, and splits. Performance is the mean of the five run-level accuracies, so it describes an individual training procedure averaged over initialization.

\subsection{Linear probe}
Multinomial logistic regression learns a linear decision rule from the original feature coordinates. For training patch \(j\), let \(z_j\in\mathbb{R}^{2560}\) denote its feature vector and \(y_j\) its class. With \(N\) training patches and \(K\) classes, the weights \(W\in\mathbb{R}^{K\times2560}\) and intercept \(b\in\mathbb{R}^{K}\) minimize
\begin{equation}
\mathcal{L}_{\lambda}(W,b)
= -\frac{1}{N}\sum_{j=1}^{N}
\log\left[\operatorname{softmax}(Wz_j+b)\right]_{y_j}
+\frac{\lambda}{2}\lVert W\rVert_F^2.
\label{eq:linear}
\end{equation}
\noindent The first term is the mean cross-entropy loss; the second penalizes large weights, with strength \(\lambda\), using the squared Frobenius norm \(\lVert W\rVert_F^2\), the sum of squared weight entries. The intercept is unpenalized. Features retain their original scale, and training patches receive equal weight. The regularization grid is
\begin{equation}
\lambda\in\{10^{-6},10^{-5},10^{-4},10^{-3},10^{-2},10^{-1},1,10,10^{2}\}.
\label{eq:lambda}
\end{equation}
\noindent A deterministic full-batch solver fits the convex objective to numerical convergence. Prediction selects the class with the largest fitted logit. This readout tests whether a regularized linear rule generalizes better than the original MLP.

\subsection{Nearest neighbours}
Cosine \(k\)-nearest neighbours (\(k\)-NN) predicts a patch's class from the labels of similar training patches. Each feature is normalized to unit length, \(u=z/\lVert z\rVert_2\). Similarity to training patch \(j\) is then the inner product \(u^{\mathsf T}u_j\). If \(\mathcal{N}_k(u)\) denotes the \(k\) most similar training patches, prediction uses an equal vote:
\begin{equation}
\widehat y(u)=\arg\max_{c\in\{1,\ldots,K\}}
\sum_{j\in\mathcal{N}_k(u)}\mathbf{1}\{y_j=c\},
\qquad k\in\{1,5,20,100,250,500\}.
\label{eq:knn}
\end{equation}
\noindent Here \(\widehat y(u)\) is the predicted class, and the indicator \(\mathbf{1}\{y_j=c\}\) equals one when neighbour \(j\) belongs to class \(c\) and zero otherwise. Neighbours are found by exact cosine search. Equal similarities are ordered by patch identifier; class ties for both probes follow a fixed class order. Several neighbours can belong to the same patient, so the vote reflects local patch density.

For each alternative readout, one hyperparameter is selected separately for each dataset and support condition by maximizing validation patient-macro balanced accuracy, averaged equally across the three splits. Numerical ties favour larger \(\lambda\) or \(k\). Test predictions use the selected setting with each split's training-only fit or neighbour bank. The MLP retains its historical hyperparameter search, so gains compare complete fitted classifiers, including their regularization and selection procedures.
\section{Endpoints}
\label{sec:analysis}
Evaluation asks three questions: does another classifier improve on the MLP under shortage, how much does broader patient coverage help, and is the alternative classifier's advantage especially large under shortage? All three use the same accuracy measure, which gives equal importance to patients within each class and then to the classes.

\subsection{Accuracy on unseen patients}
The primary endpoint is \emph{patient-macro balanced accuracy}. Its calculation has three steps. First, for each patient and class, the fraction of correctly classified patches is calculated. Second, these fractions are averaged across patients with that class. Third, the class averages receive equal weight. Thus a patient with many patches cannot dominate a class's score, and a common class cannot dominate the overall score.

For split \(r\), class \(c\), and patient \(i\), let \(\mathcal{J}_{irc}\) be the set of evaluated patches from that patient whose true class is \(c\). The patient's recall for that class is
\begin{equation}
q_{irc}=\frac{1}{|\mathcal{J}_{irc}|}
\sum_{j\in\mathcal{J}_{irc}}\mathbf{1}\{\widehat y_j=c\}.
\label{eq:patient-recall}
\end{equation}
\noindent Here \(|\mathcal{J}_{irc}|\) is the number of those patches, \(\widehat y_j\) is the predicted class of patch \(j\), and the indicator counts correct predictions. Let \(\mathcal{P}_{rc}\) be the patients with at least one evaluated patch of class \(c\). Their recalls are averaged within each class and then across the \(K\) classes:
\begin{equation}
A_r=\frac{1}{K}\sum_{c=1}^{K}
\left(\frac{1}{|\mathcal{P}_{rc}|}\sum_{i\in\mathcal{P}_{rc}}q_{irc}\right).
\label{eq:ba}
\end{equation}
\noindent The inner mean gives each eligible patient equal weight; the outer mean gives each class weight \(1/K\). A patient can contribute to several classes, once within each class. Validation selection applies the same measure to validation patients. Secondary outcomes are the individual class recalls and patch-micro balanced accuracy, which pools patches within each class before averaging class recalls.

\subsection{Classifier gain, coverage benefit, and interaction}
For one dataset, \(A_{h,s}\) denotes patient-macro balanced accuracy for classifier \(h\) under support condition \(s\), averaged equally over the three test splits. The conditions are concentrated (\(\conc\)) and spread (\(\spr\)); the MLP's five initialization-level accuracies are averaged within each split first. The three contrasts are
\begin{align}
G_{h,s}&=A_{h,s}-A_{\mathrm{MLP},s},\label{eq:gain}\\
B_h&=A_{h,\spr}-A_{h,\conc},\label{eq:breadth}\\
I_h&=G_{h,\conc}-G_{h,\spr}=B_{\mathrm{MLP}}-B_h.
\label{eq:interaction}
\end{align}
\noindent Classifier gain \(G_{h,s}\) compares an alternative readout with the MLP while holding support fixed. The central contrast, \(G_{h,\conc}\), asks whether the shortage data support better predictions than the MLP achieves. Coverage benefit \(B_h\) holds the classifier fixed and measures the improvement from broader patient coverage. Interaction \(I_h\) compares the classifier gains across support conditions: a positive value means that the alternative's advantage is larger under shortage. All differences are expressed in percentage points, obtained by multiplying differences in accuracy proportions by 100.

\begin{table}[htbp]
\centering
\small
\renewcommand{\arraystretch}{1.18}
\begin{tabularx}{\textwidth}{@{}p{0.20\textwidth}XX@{}}
\toprule
Contrast & Question & Illustrative calculation \\
\midrule
Classifier gain \(G_{h,\conc}\) & Does another classifier use the shortage data better? & Under concentrated support, MLP accuracy of \(63\%\) and linear accuracy of \(67\%\) give a gain of \(4\) points. \\
Coverage benefit \(B_h\) & How much does this classifier gain from more patients? & Linear accuracy of \(65\%\) under concentrated and \(70\%\) under spread support gives a benefit of \(5\) points. \\
Interaction \(I_h\) & Is the classifier advantage larger under shortage? & A gain over the MLP of \(5\) points under concentrated support and \(1\) point under spread support gives an interaction of \(4\) points. \\
\bottomrule
\end{tabularx}
\caption{The three contrasts separate classifier choice, patient coverage, and their interaction. Numbers are independent hypothetical examples, not experimental results.}
\label{tab:contrasts}
\end{table}

\noindent A positive gain shows usable predictive structure under the same shortage allocation. Positive coverage benefits across all readouts support a role for broader patient coverage. An interaction needs both absolute accuracies and classifier gains for interpretation: a poorly performing probe can have a small coverage benefit, and therefore a positive interaction, without being a useful alternative.

\subsection{Uncertainty across patients and initializations}
Patches from the same patient share biological and acquisition characteristics, so their prediction errors can be correlated. Treating every patch as an independent observation would overstate the amount of independent evidence. Uncertainty is therefore assessed at the patient level.

A \emph{bootstrap} approximates sampling variation by repeatedly resampling or reweighting the observed data and recalculating the statistic of interest. The \emph{Bayesian bootstrap} assigns random positive weights to observations instead of drawing a new sample with replacement. In each of 2,000 replicates here, the weights for the distinct test patients within a dataset are drawn jointly from a \(\mathrm{Dirichlet}(1,\ldots,1)\) distribution. These weights sum to one and treat patients symmetrically before the draw.

Each patient's weight is shared across classifiers, support conditions, classes, and appearances in different test splits, following the benchmark's paired analysis \citep{queisler2026benchmarkProtocol}. The within-class patient mean in Equation~\eqref{eq:ba} becomes a weighted mean, normalized by the total weight of patients present in that class and split. Class and split averages remain equally weighted. Every contrast is then recalculated using the same patient weights on both sides. This pairing preserves the correlation between predictions for the same patients when estimating uncertainty in their differences.

Each replicate also samples five initialization indices with replacement from the MLP's five confirmation runs. The same draw is shared across support conditions, splits, and contrasts involving the MLP; deterministic probe predictions remain fixed. The 2.5th and 97.5th percentiles of the resulting distributions define the reported 95\% intervals. BRACS and TCGA-UT are evaluated separately, with split-level estimates complementing the pooled estimates.
\section{Interpretation}
\label{sec:interpretation}
One percentage point is the prespecified practical-relevance threshold for the three signed contrasts. A point estimate above one point is practically promising; a lower interval bound above one point supports an improvement exceeding that threshold. A positive estimate with an interval crossing zero remains uncertain. Conversely, an interval containing zero does not establish equivalence or absence of useful information.

\begin{table}[htbp]
\centering
\small
\renewcommand{\arraystretch}{1.18}
\begin{tabularx}{\textwidth}{@{}p{0.29\textwidth}X@{}}
\toprule
Observed pattern & Supported interpretation and limit \\
\midrule
Positive shortage gain \(G_{h,\conc}\) & The alternative readout exploits predictive information that the original MLP pipeline leaves unused under the same shortage allocation. It does not identify whether optimization, regularization, or decision geometry explains the gain. \\
Positive \(G_{h,\conc}\) and \(I_h\) & The gain is larger under shortage than under spread support. If both exceed the practical threshold with precise intervals, the readout warrants a later shortage-focused study. \\
Positive \(B_h\) across readouts & Broader training coverage benefits all tested approaches. If neither probe improves under shortage, a coverage limitation remains the leading explanation within this diagnostic, not proof of intrinsic inseparability. \\
Probe gain and residual \(B_h>0\) & Both explanations can contribute: the classifier can improve while broader patient coverage still adds useful evidence. \\
Weak, conflicting, or imprecise effects & Absolute accuracies, intervals, and class recalls describe the uncertainty; neither explanation is clearly favoured. \\
\bottomrule
\end{tabularx}
\caption{Interpretation requires effect magnitude and uncertainty. None of these outcomes establishes a Bayes-error bound or an oracle ceiling.}
\label{tab:interpretation}
\end{table}

\noindent A linear-probe gain shows that a simpler linear decision rule generalizes better under the tested fitting procedure. A nearest-neighbour gain shows usable local structure in the chosen metric. Failure of either probe is weaker evidence: linear rules may miss nonlinear boundaries, while cosine neighbourhoods may reflect patient identity or unsuitable feature scaling. Failure of all three readouts cannot prove that no classifier can separate the classes or that new representations are necessary.

Broad-support success supplies a useful positive control: the same frozen representation supports the observed performance when labelled patient coverage is larger. A shortage-trained probe approaching the spread-trained MLP suggests that classifier choice can offset part of the observed deficit. Its direct difference and uncertainty quantify that comparison, while any further improvement of the spread-trained probe identifies a remaining coverage benefit.

\section{Limitations}
\label{sec:budget}
The comparison concerns fixed allocations and cohorts whose benchmark outcomes motivated the question. The intervals describe held-out patient and MLP initialization variation conditional on these allocations, selected hyperparameters, and existing splits. They omit variation from new training-patient samples and hyperparameter selection. Repeated splits share patients, and paired bootstrap weights only account for part of the dependence: overlapping training and validation data remain a limitation. Fresh patient draws or an independent held-out cohort would strengthen confirmation.

The MLP and alternative readouts differ in their fitting and hyperparameter searches as well as their decision rules. A gain therefore identifies a better-performing classification procedure; attributing it specifically to architecture, optimization, or regularization requires further comparisons. All contrasts are exploratory, and the intervals describe individual effects without simultaneous coverage across comparisons.

The findings concern patch classification with one frozen encoder and the tested patient-coverage interventions. A useful shortage-trained probe establishes available predictive structure, while persistent spread benefits establish value from broader labelled coverage for the tested readouts. Failure to find a gain leaves open the possibility that other classifiers or representations could perform better.
\bibliographystyle{plainnat}
\bibliography{references}
\end{document}
